% Generated draft: requires mathematical and source review.
\documentclass[11pt]{article}
\usepackage[margin=25mm]{geometry}
\usepackage{fontspec}
\setmainfont{Latin Modern Roman}
\usepackage{amsmath,amssymb,mathtools,unicode-math}
\setmathfont{Latin Modern Math}
\usepackage{enumitem,needspace,xurl,hyperref}
\setlength{\parindent}{0pt}
\setlength{\parskip}{5pt}
\setlength{\emergencystretch}{3em}
\newcommand{\N}{\mathbb{N}}
\newcommand{\Z}{\mathbb{Z}}
\newcommand{\Q}{\mathbb{Q}}
\newcommand{\R}{\mathbb{R}}
\newcommand{\C}{\mathbb{C}}
\newcommand{\F}{\mathbb{F}}
\newcommand{\A}{\mathbb{A}}
\newcommand{\PP}{\mathbb{P}}
\newcommand{\E}{\mathbb{E}}
\newcommand{\eps}{\varepsilon}
\newcommand{\dd}{\,\mathrm{d}}
\begin{document}
\section*{UnsolvedMath ID: 30003286}
\textbf{Source problem code:} catalogue entry
\subsection{Definitions and mathematical statement}
Write \(m(\mathrm dx)=\rho(x)\,\mathrm dx\) for the measure appearing in the supplied record, and let \(L^2(\R^d,m)\) denote the space of square-integrable functions modulo equality \(m\)-almost everywhere. The record specifies a Dirichlet form \((\mathcal E,W)\) on this space, calls \((\mathcal E,H)\) its closed realization, and assumes that a regular representation of \((\mathcal E,W)\) is given. Denote that representation by \((\widetilde{\mathcal E},\widetilde W)\) on \(L^2(S,\mu)\); here \(S\), \(\mu\), and the correspondence with the original form are placeholders for representation data that the record does not describe.

For clarity, association of a Markov process \(X\) with a form means agreement of its transition semigroup with the semigroup \((T_t)_{t\geq0}\) determined by that form: for every bounded measurable \(f\) belonging to the relevant \(L^2\) space and every \(t>0\),
\[
P_t f(x):=\mathbb E_x[f(X_t)]=T_t f(x)
\quad\text{for almost every }x
\]
with respect to the relevant reference measure. If the process has a cemetery state \(\Delta\), use the convention \(f(\Delta)=0\). This definition does not supply the hypotheses needed for a particular process realization or its uniqueness.

\textbf{Problem.} For the given regular representation, identify an associated stochastic process \(\widetilde X\) on its representation state space and describe its relationship, through the representation data, to a process \(X^H\) associated with \((\mathcal E,H)\) on the original space, whenever these process realizations exist.

\textbf{Unresolved qualifications.} The record does not define \(W\) or \(H\), explain the phrase ``closed realization,'' specify the relationship between these domains, or give the regular representation and its correspondence maps. Under the usual definition, a Dirichlet form is already closed, so the phrase alone does not determine a second domain \(H\). The record also omits conditions on \(\rho\), regularity or other hypotheses ensuring process realizations, and any boundary or killing conditions. Consequently, this is a conditional formulation of the supplied identification problem; the record does not determine a unique concrete process or a precise comparison theorem.
\subsection{Short English statement}
Given a Dirichlet form and a chosen regular representation, identify the stochastic process represented by that form and explain its relationship to the process associated with the stated closed realization \((\mathcal E,H)\). The domains, representation maps, and process-existence hypotheses require clarification before a fully specified mathematical target can be stated.
\subsection{Sources}
\begin{itemize}
\item Supplied local UnsolvedMath registry record, permanent ID \(30003286\), source problem code \texttt{catalogue entry}, entitled ``Processes Associated with Regular Representations of Dirichlet Forms.'' No individual proposer is supplied; the proposed year is recorded as 2016.
\item Supplied report citation: \emph{Heat Kernels, Stochastic Processes and Functional Inequalities} (2017), Oberwolfach Reports, pp.~3087--3147, DOI \texttt{10.4171/owr/2016/55}. \url{https://doi.org/10.4171/owr/2016/55}. The background describes Oberwolfach Reports as records of mathematical workshops and their open-problem discussions.
\item Supplied literature source: Li (2025). \url{https://doi.org/10.2969/jmsj/92649264}. No title or more specific attribution is supplied.
\item Supplied literature source: Oberwolfach Report 55 (2016). \url{https://ems.press/content/serial-article-files/46659?nt=1}.
\item Literature assessment dated exactly \texttt{2026-08-21}: the supplied curation labels the problem open and states that regular representations exist and concrete examples are understood, while a general description of the associated process and its relationship to the original form remains unresolved. This is supplied dated curation derived from the OWR material, not an independently verified literature assessment. Its time-sensitive status and cited claims require independent verification.
\end{itemize}
\end{document}
